\documentclass[11pt]{amsart}

\usepackage{fontspec}
\setmainfont{TeX Gyre Pagella}

\usepackage{amsmath,amssymb,amsthm}
\usepackage{mathtools}
\usepackage[margin=1.1in]{geometry}
\usepackage{hyperref}

\theoremstyle{plain}
\newtheorem{theorem}{Theorem}[section]
\newtheorem{proposition}[theorem]{Proposition}
\newtheorem{lemma}[theorem]{Lemma}
\newtheorem{corollary}[theorem]{Corollary}
\newtheorem{conjecture}[theorem]{Conjecture}

\theoremstyle{definition}
\newtheorem{definition}[theorem]{Definition}
\newtheorem{example}[theorem]{Example}

\theoremstyle{remark}
\newtheorem{remark}[theorem]{Remark}

\makeatletter
\let\MakeUppercase\@firstofone
\makeatother

\hypersetup{
  colorlinks=true,
  linkcolor=black,
  urlcolor=black,
  pdftitle={Retention Before Retrieval},
  pdfauthor={Flyxion}
}

\title{Retention Before Retrieval: \\
Sparse Distinction and the Cost of Carrying a Past}
\author{Flyxion}
\address{Independent Researcher}
\date{\today}

\begin{document}
\maketitle

\begin{abstract}
A system that can recover its past and a system whose past has changed
it are not the same kind of system, even when their outputs coincide.
This essay draws that distinction precisely, under the names retrieval
and retention, and argues that a large class of memory architectures
--- including the context window of a language model --- are retrieval
systems in a sense that makes them structurally unsuited to carrying
identity, whatever else they are suited to. What retention actually
requires is not more storage but a different object to store: not the
event, but the distinction the event produced. This reframes memory as
a resource-allocation problem rather than a storage problem, and
connects to a second, independently motivated idea from biological
sparsity: that sparse retention is not merely an engineering
optimization but a consequence of admissible resource allocation under
energetic, informational, and temporal constraint, in which maintaining
a distinction costs something and only some distinctions earn that
cost back. The essay closes by stating what actually survives this
argument: a retained context should scale with a signal's demonstrated
contribution to future navigation, not with the raw volume of what was
observed.
\end{abstract}

\newpage
\section{Two Questions a Memory System Can Answer}

Ask a memory system what it can do, and two different questions get
asked interchangeably, though they are not the same question. The
first is whether the past can be recovered: given the right query, can
the system reproduce or reconstruct something that happened before. The
second is whether the past has changed what the system is: independent
of whether any particular fact can be reproduced on demand, has prior
experience altered the structure that is now doing the answering. A
notebook answers the first question well and the second not at all. A
habit answers the second question well and often cannot answer the
first --- ask someone why their hands know a keyboard layout and they
frequently cannot recover the history that produced the knowledge, only
exhibit it. A database retrieves. A personality retains. The two
capacities are logically independent, and a system can have either
without the other.

\begin{definition}
A system exhibits \emph{retrieval} with respect to some past experience
if, given an appropriate query, it can reproduce or reconstruct
information about that experience. A system exhibits \emph{retention}
with respect to some past experience if that experience has altered the
system's present structure in a way that persists independent of
whether the experience itself can be reproduced.
\end{definition}

These are not opposites and not points on a single scale. A system can
retrieve without retaining, as a notebook does: the pages do not change
because they were read, and nothing about the notebook's future
behavior differs for having been consulted. A system can retain without
retrieving, as procedural memory does: the capacity persists and shapes
behavior while the originating episodes may be entirely inaccessible to
recall. And a system can do both or neither. What is being separated
here is not depth of storage but kind of relationship to the past: one
in which the past is an object the system can point to, and one in
which the past is a cause the system cannot fully separate itself from.

\begin{remark}
Large-context language models are, by this definition, sophisticated
retrieval systems. A longer context window increases what can be
reproduced on demand; it does not, on its own, alter the weights doing
the reproducing. Whatever continuity such a system exhibits across a
long conversation is continuity of retrieval, not evidence of
retention, and the two should not be conflated merely because a large
enough retrieved history can make retrieval behaviorally indistinguishable
from retention over any given exchange.
\end{remark}

\section{What Compression Loses, Layer by Layer}

A natural response to the cost of retrieval at scale is compression:
rather than carrying an entire history forward, carry a smaller
representation of it. This is a reasonable engineering instinct, and it
runs immediately into a question that is easy to state and hard to
answer: compressed with respect to what is being preserved.

Consider a sequence of compressions applied to an experience. Compress
the experience into its facts, and the reasoning path that produced
those facts may be lost --- what remains is the conclusion, not the
route. Compress the reasoning path into its structure, and the
emotional or motivational significance that made the path worth taking
may be lost --- what remains is a derivation, not why anyone cared to
derive it. Compress the significance into a scalar weight or tag, and
the capacity to explain, on later inspection, why a given conclusion
was reached at all may be lost --- what remains is a number, not an
account. Each layer of compression is individually reasonable and each
one discards something that the previous layer still carried.

\begin{proposition}
\label{prop:compression-loss}
For any compression of an experience into a smaller representation,
there exists some later query for which the compressed representation
is insufficient, unless the compression was chosen with knowledge of
that query in advance.
\end{proposition}

This is close to a triviality once stated --- any lossy compression
loses something, by definition of lossy --- but its consequence is not
trivial. It means that a compression scheme optimized for the queries
one currently anticipates will systematically fail the queries one does
not yet anticipate, and a system whose relationship to its own past is
mediated entirely by such a scheme has, in effect, decided in advance
which parts of its history are allowed to matter later.

\section{The Event Becomes Expendable Once the Change Survives}

The alternative to compressing facts is to change what counts as the
unit worth preserving in the first place. Rather than asking what facts
should survive compression, ask what distinctions should survive it. A
record of the form \emph{on a given date, a particular paper was read}
is a fact about an event. A record of the form \emph{after that date, a
particular pair of concepts began to be treated as distinct where they
previously were not} is a record of a distinction the event produced.
The second kind of record does not require the first to remain useful.
Once a distinction has been incorporated into how a system reasons, the
originating event can be discarded without discarding what the event
was for.

\begin{definition}
Given an experience $e$ occurring within a system $S$, the
\emph{retained residue} of $e$ is the set of distinctions $S$ is able to
draw after $e$ that it was not able to draw before $e$, independent of
whether $e$ itself remains recoverable.
\end{definition}

\begin{proposition}
\label{prop:residue-sufficiency}
A system that preserves the retained residue of its experiences without
preserving the experiences themselves can exhibit full retention with
respect to those experiences while exhibiting no retrieval with respect
to them.
\end{proposition}

This is the sense in which personality is closer to retention than to
retrieval, and closer to a record of distinctions than a record of
events. Nobody carries a complete transcript of their life; what
persists is what those events changed about how one distinguishes,
evaluates, and acts, and the events themselves become expendable
exactly to the degree that their residue has already been absorbed into
the structure doing the acting. History survives by becoming structure,
not by remaining stored.

\section{Convergent States, Divergent Histories}

\begin{quotation}
\emph{Three different stories leading to the same place / Same final
state, same expression on the face / But the reasons ain't the same and
the reasons matter most / History's the thing that disappears first
like a ghost.}
\end{quotation}

The argument so far has established that a retained residue can survive
independent of the event that produced it. It has not yet addressed a
sharper problem: two or more genuinely different histories can produce
the same present state, and nothing about that state, inspected on its
own, reveals which history actually occurred.

\begin{definition}
A state $s$ of a system is \emph{history-ambiguous} if there exist at
least two distinct histories $h_1 \neq h_2$ --- sequences of events and
repairs in the sense of Section~9 --- that both terminate in $s$.
\end{definition}

\begin{proposition}
\label{prop:history-ambiguity}
If a system's retained residue is computed solely from its current
state, without reference to the sequence of events and repairs that
produced it, then any two histories converging on a history-ambiguous
state necessarily produce identical residues, whether or not those
histories were equivalent with respect to future navigation.
\end{proposition}

This sharpens a claim that has been implicit since Section~1: history
matters independent of state. A residue computed only from where a
system currently stands cannot distinguish two different roads to that
same standing, even where the roads differ in exactly the respect
Proposition~\ref{prop:allocation} says should govern future behavior.
Three sequences of repair arriving at the same present configuration
may still warrant different responses to the same future perturbation,
if what future navigation depends on is not the distinction's mere
presence but the reasons it was retained --- and state convergence, by
definition, erases those reasons from anything that inspects only the
resulting state.

\begin{definition}
A residue is \emph{path-distinguishing} if it retains enough of the
sequence of events and repairs that produced it to remain
distinguishable from the residue of any other history converging on the
same state, for at least the range of future perturbations the system
needs to navigate.
\end{definition}

A system rarely achieves this by keeping an explicit log of its own
history; Section~2 already established the cost of trying to keep
events explicitly. More often, a path-distinguishing residue survives
as a visible seam: a place where the present structure still shows the
mark of the specific choice that produced it, legible to an inspector
even where the system keeps no separate record of the choice itself.
The way a repair operator resolved a conflict --- by narrowing an
assertion's scope, by replacing it outright, or by removing it, in the
sense of Section~9's definition --- leaves a different seam in each
case, even where all three repairs arrive at the same immediate
assertion, and a path-distinguishing residue is exactly the retention of
enough of that seam to tell the three repairs apart later, should a
future perturbation depend on which of them actually occurred.

A memory architecture that discards path information whenever states
converge is cheaper than one that does not, and the cost of that
cheapness is deferred rather than avoided. A future perturbation that
would have been correctly navigated by a system that retained which
road it took will instead be navigated, incorrectly or by luck, by a
system that knows only where it currently stands. This is the same
tradeoff described in Section~5 under a different name, now stated as a
structural property of state rather than a property of any particular
experience: state convergence is a compression, exactly in the sense of
Proposition~\ref{prop:compression-loss}, and it is exactly as lossy ---
only here what is lost is not a fact but the capacity to tell two true
histories apart.

\section{The Problem of Unknown Future Value}

Proposition~\ref{prop:residue-sufficiency} makes retention sound like a
strictly better deal than retrieval: keep the residue, discard the
event, save the storage. The difficulty is that the retained residue of
an experience is frequently not knowable at the moment the experience
occurs. A conversation that appears throwaway can become, months or
years later, the origin of a sustained line of work. A note recorded
for no evident reason can turn out, on later inspection, to contain the
seed of a distinction that only became load-bearing once a much larger
structure had grown up around it. At the time of recording, a system
has no reliable way to know which of its current experiences belong to
this category and which genuinely do not matter.

This produces a real tension rather than a design choice with an
obvious answer. Retaining only what currently seems to matter risks
discarding exactly the seeds whose value was not yet visible. Retaining
everything avoids that risk at the cost of drowning in a history most
of which genuinely will never matter again. Neither extreme is
workable at any scale, and the interesting design space sits between
them, in schemes that keep cheap, low-resolution traces of everything
and progressively commit to higher-resolution retention only for what
repetition, consequence, or later recontextualization demonstrates was
worth it.

\begin{remark}
This is also why a purely retrieval-based system cannot solve the
problem by simply keeping more. Keeping more addresses the risk of
discarding a future seed only if the seed can still be found later,
which shifts the burden onto search rather than removing it, and does
nothing for the deeper claim of Section~1: a system that can find an
old seed again has not thereby been changed by it.
\end{remark}

\section{Sparsity as Admissible Resource Allocation}

A structurally similar problem appears in biological nervous systems,
and it is worth drawing out because the two problems illuminate each
other. Sparse coding --- the observation that neural populations
typically represent information using a small, changing subset of
active units rather than dense, uniform activity --- is well supported
across sensory systems, and efficient-coding accounts of this
phenomenon typically explain it as a consequence of representing
structured input with minimal redundancy. A complementary account
notes that dense neural activity is metabolically expensive: firing
consumes energy, synapses require maintenance, long-range connections
require material, and signaling pathways compete for limited resources.
Evolution therefore operates in a landscape where dense representations
are frequently more costly than sparse ones, independent of any
information-theoretic argument for sparsity on its own terms.

It matters to state the underlying physics correctly here. Biological
systems do not defy the general thermodynamic tendency toward increasing
entropy; they maintain local, low-entropy structure by exporting
entropy into their environment, sustained by a persistent free-energy
gradient. The correct claim is not that some thermodynamic force drives
systems toward low entropy, but that a free-energy gradient permits the
maintenance of localized low-entropy structure, and that maintaining
such structure is energetically favorable only under specific,
persistently supplied conditions. Sparse neural organization is one
instance of a structure maintained this way: cheap to sustain relative
to the alternative, not a violation of the underlying tendency toward
disorder but a locally purchased exception to it.

\begin{proposition}
\label{prop:sparse-tradeoff}
Under a fixed energetic, informational, and temporal budget, a system
faced with representing more structure than its budget permits is
driven toward representations that preserve the distinctions with the
highest ratio of future utility to maintenance cost, independent of any
requirement that sparsity itself be the optimized quantity.
\end{proposition}

This reframes sparsity as a consequence rather than a target. A system
is not necessarily optimizing for sparse representations as such; it is
optimizing for some combination of prediction accuracy, response speed,
robustness, and developmental simplicity under a resource constraint,
and sparse representations are frequently what that combined
optimization lands on, because a dense representation of everything
available is rarely affordable and rarely the best use of a fixed
budget even when it is affordable. The distinctive claim worth
extracting from this is not about neural sparsity specifically. It is
that sparsity, wherever it appears under resource constraint, is a
symptom of a deeper allocation problem: maintaining any distinction
costs something, and only distinctions that pay that cost back in
future utility are worth maintaining.

\section{Context as an Allocation Problem, Not a Storage Problem}

Bringing Sections~5 and~6 together produces a specific proposal about
what a memory architecture concerned with retention, rather than
retrieval, should actually track. Biological systems appear to operate
close to a default of forgetting: most sensory events leave no lasting
trace, and what does persist is weighted by salience, prediction error,
reward, novelty, and repetition, not by raw exposure. A bat does not
retain every echo it has ever processed; a human does not retain every
photon that struck the retina. What survives is a sparse set of
distinctions that repeatedly demonstrated consequence for future
behavior, and the sparsity is not a limitation imposed on an otherwise
complete memory --- it is closer to the mechanism by which a usable
memory becomes possible at all.

\begin{definition}
Given a sensing or reasoning channel $c$ that contributes to a system's
decisions over time, its \emph{demonstrated contribution} is a measure
of how often, and how strongly, information from $c$ has actually
altered the system's decisions or predictions, as opposed to how much
information $c$ has merely supplied. Throughout, \emph{future
navigation} denotes the general capacity this contribution is measured
against: successful future action, prediction, inference, or decision,
whatever form that takes for the system in question. The term is left
generic deliberately, since what counts as navigation differs between a
notebook, an organism, and an institution, but every instance below
uses it in this same sense --- a channel's standing is earned by
altering outcomes that matter to the system's continuing to function,
not by the channel's mere presence in the system's history.
\end{definition}

\begin{proposition}
\label{prop:allocation}
A memory architecture oriented toward retention rather than retrieval
should allocate retained context in proportion to a channel's
demonstrated contribution to future navigation, rather than in
proportion to the volume of what that channel has observed.
\end{proposition}

Stated as a rough proportionality rather than a precise law, this is:
retained context should scale with the product of a signal's importance
and its demonstrated future utility, not with the raw sum of everything
observed. This is close to the opposite of how many current
retrieval-oriented memory proposals for language models are built,
which typically start from the assumption that more available context
is better and then treat compression as damage control applied
afterward. A system built on Proposition~\ref{prop:allocation} instead
starts from the assumption that most experience should leave no lasting
trace by default, and that retention is something a signal has to earn
through demonstrated consequence, not something granted by default and
later pared back.

This reframes the original question. The question is no longer
\emph{can this be saved}, which is a storage question with a
permissive default answer of yes. The question becomes \emph{is this
distinction worth carrying forward}, which is an allocation question
with a restrictive default answer of no, overridden only by
demonstrated contribution. A system asking the second question is
asking, continuously, which parts of its own experience have earned the
right to become structure, and treating everything else as
appropriately expendable --- not lost, in any sense that should trouble
the system, since nothing of consequence was discarded, only
information that never became a distinction worth keeping in the first
place.

\begin{example}
\label{ex:mic-array}
A physical sensing network makes the distinction between demonstrated
contribution and observed volume concrete outside any cognitive
setting. Suppose several microphones are placed at fixed points
throughout a structure and calibrated using a known reference signal
--- a clicker sounded at marked locations --- so that each channel's
recorded arrival times can be checked against the geometry that
produced them. Calibration, here, is nothing more than an explicit
version of Proposition~\ref{prop:allocation}: a channel earns standing
in the localization system to the extent that its readings reliably
correlate with the known signal, not to the extent that it carries a
large volume of data. Two channels can receive identical sound
pressure over a day and earn entirely different standing, if only one
of them tracks the calibration signal with the timing precision the
task requires.

The same hardware is simultaneously exposed to two unrelated sources of
high-volume, low-contribution signal: ambient radio-frequency
interference in the 2.4 GHz band, and low-frequency mechanical noise
--- traffic, machinery, structure-borne rumble --- conducted through
the building itself. Both are often larger in raw volume than the
calibration signal they compete with, and neither correlates with it.
Proposition~\ref{prop:allocation} gives a uniform criterion for
treating both as noise regardless of their volume, but it does not
supply a uniform remedy, since the two sources are physically distinct
phenomena: electromagnetic interference is addressed by shielding
against radiation, while mechanically conducted noise is addressed by
decoupling and mass, and treating either as though it were an instance
of the other --- shielding electromagnetically against a mechanically
conducted signal, for instance --- accomplishes nothing. The allocation
criterion says which channels deserve retained standing; it is silent
on the physical mechanism by which a channel that fails the criterion
is actually suppressed, and conflating the two is a category error the
criterion itself gives no warning against.
\end{example}

\section{False Retention}

Proposition~\ref{prop:allocation} states that a channel earns retained
context through demonstrated contribution to future navigation. It says
nothing about whether that contribution was, in any further sense,
correct. This omission is not an oversight to be patched but a
consequence to be stated directly: survival of a residue is not
evidence of its continued utility, only evidence that it once altered
outcomes strongly or repeatedly enough to earn standing. A distinction
can pay for its own maintenance many times over while being wrong.

\begin{definition}
Given the retained residue $R(e)$ of an experience $e$, $R(e)$ is said
to exhibit \emph{false retention} if it continues to be reinforced by
demonstrated contribution in the sense of
Proposition~\ref{prop:allocation}, while no longer improving, or while
actively degrading, the system's success at future navigation.
\end{definition}

False retention is common wherever a distinction was adaptive under the
conditions that produced it and those conditions have since changed. A
heuristic learned under scarcity persists under abundance because it
continues to feel confirmed by outcomes shaped, in part, by the
heuristic's own influence on behavior. A caution learned from a single
severe event generalizes far beyond the conditions that made it
reasonable and continues to be reinforced because avoidance rarely
produces evidence against itself. An overfit pattern, biological or
computational, is retained precisely because it demonstrably improved
performance on the data that produced it, which is no guarantee it
improves performance going forward. In each case the allocation
criterion of Section~7 is satisfied --- the residue has earned its
keep, repeatedly, by the only test available to the system at the time
--- and the residue is nonetheless a liability.

\begin{remark}
This is the sense in which retention is not self-validating.
Proposition~\ref{prop:allocation} explains why a distinction persists;
it does not thereby certify that the distinction should persist. A
memory architecture built only on demonstrated contribution will
faithfully retain its own mistakes with exactly the same mechanism that
retains its genuine discoveries, since the mechanism cannot, from
inside a single evaluation, distinguish a heuristic that is still true
from one that was true once.
\end{remark}

This is where the framework developed so far runs out of resources on
its own terms. Demonstrated contribution can explain what gets kept. It
cannot, by itself, explain what should later be un-kept. That requires
a second capacity, taken up next.

\section{Retention and Repair}

If retention alone can preserve a false residue indefinitely, a memory
architecture concerned with more than accumulation needs a mechanism
for revising or removing residues that retention has already
committed to structure. Call this mechanism repair, following its
sense elsewhere in this program: not a return to some prior clean
state, but an operation that acts on an admitted structure to bring it
back into agreement with what current experience actually supports.

\begin{definition}
A \emph{repair operator} on retained residues is a map $\rho$ taking a
residue $R(e)$, together with subsequent experience that conflicts with
it, to a revised residue $\rho(R(e))$ that resolves the conflict ---
either by narrowing the conditions under which $R(e)$ is asserted, by
replacing it outright, or by removing it from the system's structure
where no revision preserves its usefulness.
\end{definition}

Retention and repair are not the same capacity, and a system with only
one of them fails in a characteristic way. A system with retention and
no repair accumulates residues without ever revising them; every
distinction, once earned, is permanent regardless of whether later
experience contradicts it, and the result is a structure that grows
increasingly rigid, defending residues that the system's own current
experience no longer supports, precisely because nothing in the
architecture is permitted to act on that disagreement. A system with
repair and no retention has the opposite failure: it can revise
whatever it is currently holding, but nothing persists long enough
between revisions to accumulate into structure at all, and the result
is a system that never commits to anything long enough to have a
history to repair.

\begin{proposition}
\label{prop:retention-repair}
Retention without repair produces rigidity; repair without retention
produces forgetfulness. A memory architecture requires both, applied to
the same object: repair does not compete with retention but acts on
what retention has already produced.
\end{proposition}

This gives false retention, from Section~8, a resolution rather than
merely a diagnosis. A false residue is not a residue that should never
have been retained --- at the time it was retained, it satisfied every
criterion the system had available. It is a residue that has since
accumulated conflicting evidence and is now a candidate for repair. The
allocation criterion of Proposition~\ref{prop:allocation} governs what
enters structure; the repair operator of this section governs what
leaves it or is revised within it; and a healthy memory architecture is
one in which both operate continuously on the same growing set of
residues, rather than one in which retention runs and repair is called
in only after failure becomes catastrophic.

\section{Collective Retention}

Everything so far has been stated for a single system carrying its own
history forward. The same distinction between retrieval and retention
applies at a different scale, where the carrier of a residue is not one
persistent system but a population of systems, individually mortal, none
of which needs to survive for the residue to persist.

A scientific field retains distinctions --- which questions are worth
asking, which methods are trustworthy, which results require
extraordinary evidence --- that no individual scientist working within
it originated or could fully reconstruct from first principles. A
language retains grammatical and semantic distinctions across many
generations of speakers, none of whom individually invented them and
most of whom could not explain their historical origin, yet every
fluent speaker's behavior is shaped by them. A civilization preserves
techniques, from metallurgy to record-keeping, long after the specific
events and specific people that first produced them are gone.

\begin{definition}
\emph{Collective retention} occurs when a distinction persists across
the replacement of its individual carriers: no single member of the
population that carries the distinction forward needs to survive, or
even to have been present when the distinction was first produced, for
the population as a whole to continue exhibiting it.
\end{definition}

This is retention in the full sense defined in Section~1, not a
metaphorical extension of it: the population's present structure has
genuinely been altered by past experience, independent of whether any
particular past experience, or any particular past carrier of it, can
now be retrieved. What differs from individual retention is only the
substrate. An archive belonging to such a population is a retrieval
system in exactly the sense of Section~1 --- it can be consulted, and
consulting it does not itself alter the population's behavior. The
institution, language, or practice built on top of the archive is
where retention actually lives, in the same sense that a habit, not a
notebook, is where an individual's retention lives.

\begin{remark}
Section~5's problem of unknown future value appears again here, at
larger scale and with higher stakes: a distinction that looks
dispensable within one generation's frame of reference may be exactly
the seed a later generation needs, and a population that discards it
because no current member sees its value has no mechanism for
recovering the loss once the last carrier who could reconstruct it is
gone. Collective false retention, and collective repair in the sense of
Section~9, are both real and both harder to execute than their
individual counterparts, since no single carrier has the authority or
the visibility to perform either alone.
\end{remark}

\begin{example}
\label{ex:incarnations}
A system built by interchanging language-model instances across
successive incarnations --- swapping which trained model performs a
given step, delegating to auxiliary neural networks where a language
model's own competence runs out, searching a problem space in something
closer to a heuristic search than to one continuous train of thought
--- is a collective-retention system in exactly the sense above: no
single incarnation needs to persist, or even to have been present for
an earlier decision, for the system as a whole to continue acting on
distinctions earlier incarnations established. What is available inside
one incarnation's context window is retrieval, in the sense of
Section~1, and it vanishes when that incarnation is replaced. What
survives the replacement has to live somewhere else --- a search
history, a scoring function, an accumulated set of heuristics external
to any single incarnation's weights or context --- which is the
archive-versus-institution split of this section applied to a system
built rather than grown.

This sharpens a question that is often asked more loosely than it
needs to be: whether such a system understands anything it does. A
single incarnation producing a correct-sounding correlation once is an
act of retrieval and settles nothing about understanding on its own,
regardless of how fluent it looks. Understanding, on the terms
developed in this essay, would be the stronger and more specific claim
that the correlation has become retention: that some later incarnation,
with no access to the episode in which the correlation was first drawn
and no query capable of retrieving it, nonetheless acts as though it
had been. Whether that transfer actually occurs is a question about the
system's architecture --- whether it has anywhere for retention to
live outside any single incarnation's retrieval --- and not a question
that any one incarnation's performance, however impressive, can settle
by itself.
\end{example}

\section{Identity as a Compression Scheme}

Return to the individual case with Sections~8 and~9 in hand. Nobody
carries a complete transcript of their life, and Section~3 already
argued that the transcript becomes expendable once its residue
survives. What has not yet been stated is what the accumulated set of
retained residues, taken together, actually amounts to for a single
system across a long enough history. The claim of this section is that
it amounts to something worth calling identity, and that identity is
best understood as a compression scheme operating over residues rather
than a further object added on top of them.

\begin{definition}
Given a system $S$ with a history of experiences $e_1, e_2, \ldots$,
its \emph{identity at time $t$} is the compressed structure obtained
from the accumulated, repaired residues
$\rho(R(e_1)), \rho(R(e_2)), \ldots$ retained up to $t$ --- not the
residues individually enumerated, but whatever further compression of
them the system uses to act coherently without consulting each one
explicitly.
\end{definition}

This is compression in the same sense discussed in Section~2, and
Proposition~\ref{prop:compression-loss} applies to it exactly as it
applies to any other compression: an identity, so defined, will fail
some future query that a complete record of residues would have
answered. This is not a flaw particular to identity; it is the same
tradeoff every layer of this essay has been describing, now applied one
level higher, to the residues themselves rather than to the events that
produced them. What makes the compression worth performing is the same
argument as before --- carrying every residue explicitly forward is as
unworkable as carrying every event forward, and a system's sense of
\emph{the kind of system I am} is what lets it act without re-deriving
its position from first principles on every occasion.

\begin{remark}
This reframes a familiar observation. Personality changing only slowly,
and only under sustained pressure, is not evidence that personality is
a fixed thing weakly connected to experience; it is exactly what a
compression scheme operating over demonstrated contribution, in the
sense of Proposition~\ref{prop:allocation}, should be expected to do.
A single experience rarely re-earns enough demonstrated contribution to
shift a compression built from thousands of prior ones; a sustained,
repeated pattern of new experience can.
\end{remark}

\section{Retention Without Retrieval}

The argument to this point has stayed close to cognitive systems ---
individuals, institutions, languages --- where retrieval is at least
conceivable even where it does not occur. It is worth pushing the
distinction from Section~1 to a point where retrieval is not merely
absent but not conceivable at all, since that is where the distinction
stops being a claim about memory architectures specifically and becomes
a more general claim about historical influence.

A tree that survives a drought retains the drought: its growth rings,
root structure, and subsequent water regulation are altered by the
event in ways that persist for the rest of its life, and none of this
requires the tree to retrieve, in any sense, a representation of the
drought that produced it. An immune system that has encountered a
pathogen retains the encounter in the form of memory cells calibrated
to respond faster to a future exposure, with nothing resembling
retrieval of the original infection as an event. A muscle subjected to
training retains the training as altered fiber composition and
capacity, retrievable by no process the organism has access to, only
exhibited under later load. A species subjected to a selection pressure
retains it in the composition of the population generations later, with
no individual member of that population, and no archive belonging to
it, storing anything that could be called a representation of the
originating pressure.

\begin{proposition}
\label{prop:retention-without-retrieval}
Retention does not require retrieval to be available even in
principle. A system can be fully and permanently changed by a past
event while possessing no mechanism, actual or possible, for
representing that event as an object to be recovered.
\end{proposition}

This is the strongest form the distinction from Section~1 can take, and
it changes what kind of claim the essay is making. Retrieval is a
relation to a representation of the past: it requires that the past be
stored, in some form, as an object the system can later point to.
Retention is a relation to a transformation produced by the past: it
requires only that the system now be different, in some functionally
relevant way, than it would have been absent the event, whether or not
anything resembling a representation of that event exists anywhere. Once
stated this way, the essay is no longer primarily about memory
architectures, language models, or even cognition. It is a claim about
what it means for a past to have happened to a system at all --- that
having happened is a fact about present structure, not a fact about
present or future access to a stored account of it.

\section{Crystallized Claim}

Retrieval and retention answer different questions, and a system built
to excel at the first does not thereby make progress on the second: a
larger context window is more retrieval, not more retention, and the
two should not be treated as the same capacity at different scales.
What retention actually requires is a change in what counts as the
unit of memory, from the event to the distinction the event produced,
since a system that retains distinctions rather than events can carry
its history forward while discarding almost all of the history itself,
provided the distinctions have genuinely been incorporated into how the
system reasons. This does not resolve the difficulty of not knowing, at
the moment of an experience, which of its distinctions will later prove
load-bearing; that difficulty is real and is not solved by keeping
more, since keeping more only defers the problem to retrieval without
ever producing retention. What it does provide is a criterion for what
a memory system should be trying to do under that uncertainty: allocate
retained structure in proportion to a signal's demonstrated
contribution to future navigation rather than to the volume of
experience available, treating sparsity not as a compression technique
applied to an otherwise complete record but as the expected consequence
of any system honestly pricing the cost of carrying a distinction
forward.

That criterion is not self-validating, and taking it seriously forced
the argument further than a theory of storage needs to go. A residue
earns standing by demonstrated contribution, and demonstrated
contribution is not the same fact as continued correctness; a
distinction can pay for its own maintenance many times over while
actively working against the system that keeps it. Retention alone has
no way to tell the two apart, which is why retention cannot be the
whole of a memory architecture. It requires a second capacity running
alongside it, repair, that acts on retained residues rather than
competing with the act of retaining them: retention without repair
hardens into rigidity, repair without retention never accumulates into
anything a system could call a history, and a working memory
architecture is the two operating continuously on the same growing set
of residues.

Once residues, rather than events, are the object of concern, three
further consequences follow that a storage-oriented theory of memory
would have no occasion to notice. The same distinction that separates
retrieval from retention inside one system also separates an archive
from an institution across many: a population can retain what no
individual member of it could reconstruct, provided the distinction
survives the replacement of its carriers, which is retention in the
full sense of Section~1 rather than a metaphor borrowed from it. Within
a single system, the accumulated and repaired residues of a long
history are not carried forward individually but compressed further
still, into whatever a system uses to act coherently without consulting
its history explicitly --- and that further compression is what this
essay is willing to call identity, subject to the same lossy tradeoffs
as any other compression discussed here, no more and no less. And
retention, pushed to its limit, turns out not to require retrieval even
in principle: a tree, an immune system, a muscle, a population under
selection are all permanently altered by a past none of them represents
anywhere, which is the clearest available evidence that retention is a
relation to a transformation produced by the past, while retrieval is
only ever a relation to a stored representation of it, and the two
were never the same kind of fact about a system to begin with.

Memory, on this view, is not a storage problem that compression happens
to help with. It is a resource-allocation problem that storage happens
to be one input to, revision is a second and equally necessary input
to, and the thing actually being allocated is not space but standing:
which distinctions a system's past has earned the right to leave behind
in what the system currently is, subject to continuous revision as
demonstrated contribution and continued correctness come apart, at
every scale from a single organism to a civilization, whether or not
anything answerable to the name of memory is present to notice the
allocation happening at all.

\bigskip
\begin{thebibliography}{9}

\bibitem{barlow1961}
H. B. Barlow, \emph{Possible principles underlying the transformation
of sensory messages}, in Sensory Communication, MIT Press, 1961.

\bibitem{olshausen1996}
B. A. Olshausen and D. J. Field, \emph{Emergence of simple-cell
receptive field properties by learning a sparse code for natural
images}, Nature, 1996.

\bibitem{friston2010}
K. Friston, \emph{The free-energy principle: a unified brain theory?},
Nature Reviews Neuroscience, 2010.

\bibitem{tulving1985}
E. Tulving, \emph{Memory and consciousness}, Canadian Psychology, 1985.

\bibitem{baddeley2000}
A. Baddeley, \emph{The episodic buffer: a new component of working
memory?}, Trends in Cognitive Sciences, 2000.

\end{thebibliography}

\end{document}

